\documentclass[11pt]{article}

\usepackage[margin=0.85in]{geometry}
\usepackage[T1]{fontenc}
\usepackage[utf8]{inputenc}
\usepackage{amsmath,amssymb}
\usepackage{xcolor}
\usepackage{booktabs}
\usepackage{tabularx}
\usepackage{tcolorbox}
\usepackage{hyperref}

\definecolor{auditred}{RGB}{185,38,38}
\definecolor{auditblue}{RGB}{34,80,145}
\definecolor{auditgray}{RGB}{245,246,248}

\setlength{\parindent}{0pt}
\setlength{\parskip}{0.55em}

\newtcolorbox{resultbox}[1]{
  colback=auditgray,
  colframe=#1,
  boxrule=0.8pt,
  arc=1.5pt,
  left=7pt,
  right=7pt,
  top=6pt,
  bottom=6pt
}

\begin{document}

\begin{center}
{\Large\bfseries Quick sign test for the shortcut sign}\\[0.3em]
{\small Ratio test for Eqs. (29) and (32), using real values of \(z\)}
\end{center}

\begin{resultbox}{auditblue}
\textbf{Short verdict.}
The stochastic shortcut matrix, where \(\lambda=\beta(1-q)\) is moved from the
self-loop at \(u\) to the directed edge \(u\to v\), agrees with Eq. (32) only
with the corrected plus/plus sign.  The original minus/minus sign does not
match this stochastic matrix.
\end{resultbox}

I used the proposed ratio check: compute the first-passage generating
function from a full-ring shortcut propagator by
\[
\widetilde F_{n_0}(v,z)=\frac{\widetilde S_{n_0}(v,z)}
{\widetilde S_v(v,z)}.
\]
I then compared this value with Eq. (32) using both signs.

There is one important point in the test.  If Eq. (29) is used literally as
printed in the manuscript, its ratio agrees with the original minus/minus
Eq. (32).  However, the actual probability-conserving stochastic shortcut
matrix for \(u\to v\) agrees with the corrected plus/plus Eq. (32).  Thus this
test pushes the sign inconsistency one step earlier: Eq. (29), as printed, has
the opposite defect sign if it is interpreted as moving self-loop mass from
\(u\) to \(v\).

\begin{resultbox}{auditred}
\textbf{Sanity check.}
For positive real \(z<1\), a first-passage generating function has non-negative
coefficients and therefore cannot be negative.  In the \(N=10\), \(q=2/3\),
\(\beta=2/7\), \(u=5\), \(v=0\) test, the original sign gives a negative value
for \(n_0=4\), \(z=0.8\), while the stochastic matrix and the corrected sign
give a positive value.
\end{resultbox}

\section*{Error Summary}

The grid uses two symmetric cases, including the figure parameters in the manuscript
\(N=10\), \(q=2/3\), \(\beta=2/7\), \(u=5\), \(v=0\), with several starts
\(n_0\) and real values of \(z\).  Site labels here are zero-based.

\subsection*{Main sign check}

\begin{center}
\begin{tabularx}{0.98\linewidth}{@{}Xr@{}}
\toprule
Quantity compared over the grid & Maximum absolute error \\
\midrule
Original Eq. (32) vs stochastic shortcut matrix & \(3.669746\times10^{-1}\) \\
Corrected Eq. (32) vs stochastic shortcut matrix & \(2.775558\times10^{-16}\) \\
Literal Eq. (29) ratio vs original Eq. (32) & \(1.318390\times10^{-16}\) \\
Stochastic-sign Eq. (29) ratio vs corrected Eq. (32) & \(5.551115\times10^{-17}\) \\
\bottomrule
\end{tabularx}
\end{center}

\subsection*{Control checks}

\begin{center}
\begin{tabularx}{0.98\linewidth}{@{}Xr@{}}
\toprule
Quantity compared over the grid & Maximum absolute error \\
\midrule
Eq. (19) \(W_{n_0}(u,z)\) vs absorbing matrix & \(1.043610\times10^{-14}\) \\
Eq. (28) \(W_u(u,z)\) vs absorbing matrix & \(1.687539\times10^{-14}\) \\
\bottomrule
\end{tabularx}
\end{center}

\section*{Representative Values}

\begin{center}
\small
\begin{tabular}{@{}rrrrrrr@{}}
\toprule
\(n_0\) & \(z\) & stochastic matrix & literal Eq. (29) & original Eq. (32) & corrected Eq. (32) \\
\midrule
2 & 0.50 & 0.044096501039 & 0.042961443311 & 0.042961443311 & 0.044096501039 \\
2 & 0.80 & 0.197025316799 & 0.173752283343 & 0.173752283343 & 0.197025316799 \\
2 & 0.95 & 0.532769134231 & 0.394627514541 & 0.394627514541 & 0.532769134231 \\
4 & 0.80 & 0.095953324694 & -0.033502923905 & -0.033502923905 & 0.095953324694 \\
\bottomrule
\end{tabular}
\end{center}

\section*{Why This Locates Eq. (29)}

For the stochastic shortcut interpretation, the transition mass
\(\lambda=\beta(1-q)\) is removed from the \(u\to u\) self-loop and added to
the \(u\to v\) edge.  In the resolvent, this produces the sign pattern with a
\(+z\lambda\{P_u(u,z)-P_v(u,z)\}\) denominator.  The printed Eq. (29) uses the
opposite rank-one defect sign.  Therefore the printed Eq. (29) reproduces the
original Eq. (32), but both differ from the direct stochastic shortcut matrix.

\end{document}
